% Source: https://artofproblemsolving.com/community/c3h1900141_indistinguishability
% Stable identity: aops-topic-1900141-q01
\textbf{(a)} A colour pattern is fixed by choosing which $3$ of the $8$ places hold blue marbles; the other $5$ places are red. So there are $\dbinom83=\boxed{56}$ patterns.

\textbf{(b)} Now the $8$ marbles are $8$ different objects, so there are $8!=\boxed{40320}$ rows.

\textbf{(c)} Take any colour pattern from (a). The $5$ numbered red marbles can be placed in the $5$ red places in $5!$ ways, and the $3$ numbered blue marbles in the $3$ blue places in $3!$ ways. Every row of numbered marbles arises in exactly one way like this, from its own colour pattern. Hence
\[
8!=\dbinom83\cdot5!\cdot3!,\qquad 40320=56\cdot120\cdot6 .
\]
So whether objects count as distinguishable is part of the question being asked: if only colours are visible, rows that differ by swapping two red marbles look the same and are counted once, as in (a).
